\subsubsection{de Branges 逆谱与极端零知识：canonical system、谱移等价与 RH 的算子化正定跳跃范式}\label{subsubsec:xi-debranges-canonical-rhp-adelic}
本段在 \S\ref{subsubsec:xi-hardy-inner-clark-carleson} 的 Hardy/内函数化框架上继续向“谱实现/逆谱重建”推进：
当完成化行列式在临界线上具有纯相位且无奇异内因子时，它不仅是一个内函数，而且可被唯一提升为一条 canonical system（de Branges）；
并且“极端零知识”（无纤维内混合自由度）可被完全等价地译为该 canonical system 的秩退化结构。
在此基础上，谱移（Krein）公式、确定性点过程的核极限、SU(1,1)–Riemann--Hilbert 正定跳跃以及 adelic/Dirichlet 族的统一零知识拼接，可被写入同一条“内函数 \(\Rightarrow\) 可实现 \(\Rightarrow\) 临界线集中/零自由带”的可审计推理链条。

\paragraph{半平面化与无奇异内因子假设}
沿用 \S\ref{subsubsec:xi-hardy-inner-clark-carleson} 的记号：\(D(s)\) 在 \(\mathbb{H}_{1/2}\) 上解析且为 bounded type，并满足反射对称 \eqref{eq:xi-inner-reflection-symmetry}。
再假定视界零知识幺正性（纯相位）成立：
\[
|D(\tfrac12+it)|=1\qquad(\text{a.e. }t\in\mathbb{R}),
\]
并且圆盘化内函数 \(\mathfrak{D}\) 无奇异内因子（等价于 \S\ref{subsubsec:xi-hardy-inner-clark-carleson} 中的 Clark 测度几乎处处绝对连续口径）。
把变量写为 \(s=\tfrac12+iz\)（\(\Im z>0\)），并记
\[
\mathcal{D}(z):=D(\tfrac12+iz),
\qquad z\in\mathbb{C}_+.
\]

\begin{theorem}[de Branges--视界逆谱重建与“极端零知识”判别]\label{thm:xi-debranges-canonical-extreme-zk}
\leanverified{paper\_xi\_debranges\_canonical\_extreme\_zk}
在上述假设下，存在一条 canonical system
\[
J\,Y'(x,z)=z\,H(x)\,Y(x,z),\qquad x\in[0,\infty),
\]
其中
\(
J=\begin{psmallmatrix}0&-1\\[2pt]1&0\end{psmallmatrix}
\)，\(H(x)\succeq 0\) 且 \(\mathrm{tr}\,H(x)=1\) a.e.，
使得存在一个 Hermite--Biehler 型 de Branges 函数 \(E(z)\) 满足（在上半平面意义下）
\[
\boxed{\;\mathcal{D}(z)=\frac{E^\sharp(z)}{E(z)}\;},
\qquad
E^\sharp(z):=\overline{E(\overline z)}.
\]
该 canonical system 在同构意义下唯一：两条产生同一 \(\mathcal{D}\) 的系统仅至多相差一个 a.e. 可测的 \(SO(2)\) 规范变换（对 \(H\) 作共轭）。
\par
进一步地，若把“极端零知识”表述为：对几乎处处视界实例 \(y\)，模拟器给出的条件分布是可见概率单纯形的极端点（等价于纤维内无额外混合自由度），
则该条件当且仅当 canonical system 的 Hamiltonian 满足
\[
\boxed{\;\mathrm{rank}\,H(x)=1\quad\text{a.e. }x.\;}
\]
\end{theorem}
\begin{proof}[证明要点]
由定理 \ref{thm:xi-hardy-inner-characterization} 与“无奇异内因子”假设，\(\mathcal{D}\) 是上半平面上的（无奇异部分）内函数。
de Branges 理论给出内函数的表示 \(\mathcal{D}=E^\sharp/E\) 以及与 canonical system 的对应；\(\mathrm{tr}\,H=1\) 通过尺度规一化固定。
唯一性来自 canonical system 的 gauge 等价（\(SO(2)\) 共轭不改变 \(E^\sharp/E\)）。
\par
对“极端零知识 \(\Leftrightarrow\) 秩一”：在 \(2\times 2\) 情形，满足 \(\mathrm{tr}\,H(x)=1\) 的正半定矩阵族形成一个逐点凸集，其极端点恰为秩一投影；将“纤维内无混合自由度”理解为该逐点凸集中的极端化选择，即得秩一判别。证毕。
\end{proof}

\begin{proposition}[谱移–信息泄漏恒等式：Krein 公式在零知识接口上的相对熵表述]\label{prop:xi-krein-spectral-shift-equals-kl}
\leanverified{paper\_xi\_krein\_spectral\_shift\_equals\_kl}
在定理 \ref{thm:xi-debranges-canonical-extreme-zk} 的 canonical system 模型中，以挑战种子 \(\alpha\in\TT\) 参数化自伴边界条件，得到一族自伴扩张 \(S_\alpha\) 与对应谱测度 \(\sigma_\alpha\)（可取为 Clark–谱测度，推论 \ref{cor:xi-clark-measure-horizon-semantics}）。
设在每个 \(\alpha\) 下转录分布 \(P_\alpha\) 在可见 \(\sigma\)-代数上与 \(\sigma_\alpha\) 同构（零知识分解在该口径下对齐到谱测度）。
则对任意 \(\alpha,\beta\in\TT\) 有严格恒等式
\[
\boxed{\;
D_{\mathrm{KL}}(P_\alpha\Vert P_\beta)
=\int_{\mathbb{R}}\log\!\left(\frac{\dd\sigma_\alpha}{\dd\sigma_\beta}(t)\right)\dd\sigma_\alpha(t)\;}
\]
并且在 canonical system 的正则化口径下，右端亦等于一条 Krein 谱移密度 \(\xi_{\alpha,\beta}(t)\) 的积分：
\[
\boxed{\;
\int_{\mathbb{R}}\log\!\left(\frac{\dd\sigma_\alpha}{\dd\sigma_\beta}\right)\dd\sigma_\alpha
=\int_{\mathbb{R}}\xi_{\alpha,\beta}(t)\,\dd t\; }.
\]
因此，“零知识泄漏的相对熵”与“边界条件切换的谱移”在该一维缺陷指标模型中同一。
\end{proposition}
\begin{proof}[证明要点]
第一式为同构识别后的测度相对熵恒等式。
第二式由 Krein 谱移理论给出：谱测度变化的 Radon--Nikodym 对数可用扰动行列式（谱移函数）表出，并在本模型的正规化下具有密度 \(\xi_{\alpha,\beta}\) 的积分表示。证毕。
\end{proof}

\begin{proposition}[零点点过程的行列式结构：视界均匀性 \(\Rightarrow\) 正弦核极限]\label{prop:xi-debranges-dpp-sine-kernel-limit}
沿定理 \ref{thm:xi-debranges-canonical-extreme-zk} 取秩一 Hamiltonian \(H(x)=v(x)v(x)^{\mathsf T}\) a.e.，
并假设“视界均匀性”：存在常矩阵 \(H_\infty\succeq 0\) 使 \(H(x)\to H_\infty\) 于 Ces\`aro（或更强遍历平均）意义成立，且相位增长满足线性正则。
则对几乎处处种子 \(\alpha\)，扩张 \(S_\alpha\) 的谱点集合（等价于临界线上零点高度集合）可由 de Branges 再现核构造出一族行列式点过程；在局部尺度归一化后，其核收敛到正弦核，从而二点相关函数在尺度极限满足 Montgomery--Dyson 型正弦律。
\end{proposition}
\begin{proof}[证明要点]
秩一 canonical system 给出一族 de Branges 空间，其再现核决定一个自然的行列式核。
在 Ces\`aro 均匀性与线性相位正则下，局部密度稳定并允许作等密度缩放；核的局部极限退化到常系数模型，从而得到正弦核。证毕。
\end{proof}

\begin{theorem}[SU(1,1)–Riemann--Hilbert 等价：零知识幺正性将 RH 化为正定跳跃可解性]\label{thm:xi-su11-rhp-equivalence}
设由转录更新诱导一矩阵跳跃数据 \(V(t)\)（定义于临界边界 \(t\in\mathbb{R}\)），其元素由可见挑战–响应相关核给出并满足 \(\det V(t)=1\) a.e.。
若零知识幺正性成立且外因子消失（即 \(\log|D(\tfrac12+it)|=0\) a.e.，命题 \ref{prop:xi-entropy-production-equals-outer-factor}），则必有
\[
\boxed{\;V(t)\in SU(1,1)\quad\text{a.e. }t,\;}
\]
并且存在唯一归一化解 \(M(z)\)（\(z\in\mathbb{C}_+\)）满足标准 Riemann--Hilbert 问题
\[
M_+(t)=M_-(t)\,V(t),
\qquad
M(\infty)=I.
\]
在此框架下，“非平凡零点全部位于临界线”与“该 Riemann--Hilbert 问题在正定范畴中可解且其 Beals--Coifman 算子严格可逆”互为等价模板：任一离线零点事件对应 \(SU(1,1)\) 正定性破缺所导致的不可去共振极；反之，正定可解性排除一切离线共振极，从而迫使零点集中于临界线。
\end{theorem}
\begin{proof}[证明要点]
纯相位/外因子消失把跳跃数据压入保形不变的幺正对称类（\(SU(1,1)\)）。
Riemann--Hilbert 的正定可解性等价于对应奇异积分算子（Beals--Coifman）在 \(L^2\) 上可逆；离线零点事件表现为 resolvent 的不可去极点，从而与可逆性矛盾。证毕。
\end{proof}

\begin{proposition}[Maslov 指数–谱流–熵量子化：挑战参数的拓扑刚性]\label{prop:xi-maslov-spectralflow-entropy-quantization}
在定理 \ref{thm:xi-debranges-canonical-extreme-zk} 的 canonical system 中，将挑战种子 \(\alpha\in\TT\) 解释为边界 Lagrangian 子空间 \(L_\alpha\) 的参数。
令 \(N_\alpha(T)\) 为高度截断 \([0,T]\) 上的谱流计数（等价于临界线零点计数差分）。
则对任意两种子 \(\alpha,\beta\) 有整数不变量恒等式
\[
\boxed{\;
N_\alpha(T)-N_\beta(T)=\mathrm{Maslov}\bigl(L_\alpha,L_\beta;[0,T]\bigr)\in\mathbb{Z}. \;}
\]
若进一步将“不可见纤维规模/容量”用常数 \(\kappa_{\mathrm{cap}}>0\) 归一化，并取一可加熵泛函 \(S_\alpha(T)\)（例如由对数势/账本定义），则存在量子化律
\[
\boxed{\;
S_\alpha(T)-S_\beta(T)=\kappa_{\mathrm{cap}}\bigl(N_\alpha(T)-N_\beta(T)\bigr).\;}
\]
\end{proposition}
\begin{proof}[证明要点]
谱流与 Maslov 指数在一维 canonical system 的边界 Lagrangian 参数化下同型；熵泛函的可加性把“每穿越一条谱点”贡献固定增量 \(\kappa_{\mathrm{cap}}\) 的量子化读出，得到所述恒等式。证毕。
\end{proof}

\begin{proposition}[Adelic 零知识拼接原理：有限域纯权与一致模拟 \(\Rightarrow\) 整体临界线集中]\label{prop:xi-adelic-zk-gluing-critical-line}
\leanverified{paper\_xi\_adelic\_zk\_gluing\_critical\_line}
设对每个素数 \(p\) 给定局部因子 \(Z_p(u)\)（有理函数），满足：
\begin{enumerate}
  \item \(Z_p\) 具 Weil 型纯权（全部非平凡特征根模为 \(p^{-1/2}\)）；
  \item 存在跨 \(p\) 的一致模拟器，使局部转录统计在 Chinese remainder 拼接下对可见因子闭合；
  \item Euler 拼接
  \(
  D(s):=\prod_p Z_p(p^{-s})^{-1}
  \)
  具有解析延拓与反射对称 \eqref{eq:xi-inner-reflection-symmetry}。
\end{enumerate}
则一致模拟器迫使 \(D\) 在 \(\mathbb{H}_{1/2}\) 为内函数且外因子消失（定理 \ref{thm:xi-hardy-inner-characterization} 与命题 \ref{prop:xi-entropy-production-equals-outer-factor} 的拼接口径）。
若再加“无奇异内因子”条件，则 \(D\) 的全部非平凡零点只能位于临界线 \(\Re(s)=\tfrac12\)（定理 \ref{thm:xi-rh-sufficient-template-inner-clark-conservative} 的模板化推论）。
\end{proposition}

\begin{proposition}[统一零知识对 Dirichlet 族的零自由带：排除 Siegel 型例外零]\label{prop:xi-dirichlet-zero-free-from-unified-zk}
考虑 Dirichlet 角色族 \(\chi\bmod q\) 的“素数挑战视界协议”，其可见实例仅包含模 \(q\) 的残余类统计。
假设存在统一模拟器，使对所有 \(q\)：
转录长度为 \(\mathrm{polylog}(q)\)，且健全性误差至多 \(\exp(-c\log q)\)（常数 \(c>0\) 与 \(q\) 无关）。
则任一实角色 \(\chi\) 的假设性例外实零 \(\beta_\chi\)（若足够靠近 \(1\)）将诱导低复杂度区分器，从而与统一零知识不可区分性矛盾。
因此存在常数 \(c'>0\) 使得
\[
\boxed{\;\beta_\chi\ \le\ 1-\frac{c'}{\log q}\; }.
\]
\end{proposition}
\begin{proof}[证明要点]
例外零在可见素数残余类统计上诱导可放大的偏置；在 \(\mathrm{polylog}(q)\) 尺度内即可构造区分器使 TV 距离超过 \(\exp(-c\log q)\) 的健全性阈值，从而与统一模拟器矛盾。证毕。
\end{proof}

\begin{corollary}[临界零点的几乎处处简单性：绝对连续挑战种子排除多重谱交叉]\label{cor:xi-ae-simple-zeros-under-ac-seeds}
在定理 \ref{thm:xi-debranges-canonical-extreme-zk} 的扩张族 \(\{S_\alpha\}_{\alpha\in\TT}\) 中，谱点多重性等价于边界 Lagrangian 数据的非横截交叉。
若挑战种子分布在 \(\TT\) 上具有绝对连续密度，且零知识正则边界条件保证边界 Weyl 函数的 nontangential 边界值存在并满足非退化导数条件，
则产生多重谱点（多重零点）的种子集合为 Lebesgue 零测。
\end{corollary}

\begin{proposition}[非正规性障碍定理：离线零点当且仅当视界更新算子不可相似正规]\label{prop:xi-nonnormality-obstruction-offcritical}
令 \(\mathsf{T}\) 为作用于 \(L^2(Y,\nu)\) 的“视界转录更新算子”（由一步挑战–响应统计核定义），并假设 \(D(s)\) 可由 \(\mathsf{T}\) 的特征函数/行列式正规化得到且满足反射对称 \eqref{eq:xi-inner-reflection-symmetry}。
则在与参考态相容的实现口径下，
\[
\boxed{\;
\exists\ \text{离临界线零点事件}\quad\Longleftrightarrow\quad
\mathsf{T}\ \text{不可相似为正规算子}. \;}
\]
特别地，若强制详细平衡（从而获得幺正散射切片与零知识幺正性），则 \(\mathsf{T}\) 必相似正规并排除全部离线零点事件。
\end{proposition}
\begin{proof}[证明要点]
相似正规等价于在某等价内积下满足谱定理与正交分解，从而其特征函数在边界呈纯相位并无离线共振极；
反之，离线零点事件对应 resolvent 的不可去极点，必破坏任何与参考态相容的正规化。证毕。
\end{proof}

\begin{proposition}[证明长度下界与离线零点密度：BM 密度强迫信息复杂度线性增长]\label{prop:xi-bm-density-forces-linear-transcript-length}
设 \(\Re(s)>\tfrac12\) 的离线零点集合在高度变量上具有正的 Beurling--Malliavin 下密度（以可见尺度归一化）。
并假设在高度尺度 \(T\) 上的全部可操作检验均可由长度至多 \(m(T)\) 的转录实现，且零知识要求保持 \(\varepsilon(T)\)-不可区分（\(\varepsilon(T)\to 0\)）。
则必存在常数 \(c>0\) 使
\[
\boxed{\; m(T)\ \ge\ c\,T\;}
\]
从而任何 \(\mathrm{polylog}(T)\) 级别的证明长度与正密度离线零点不相容。
\end{proposition}
\begin{proof}[证明要点]
正密度离线零点对应临界线上的 Poisson 共振项以正密度叠加（命题 \ref{prop:xi-phase-derivative-poisson-superposition}），其可放大总变差偏离在尺度 \(T\) 上至少线性增长；
若转录长度不足以线性增长，则存在低复杂度区分器迫使 \(\varepsilon(T)\) 不能趋于零，与零知识假设矛盾。证毕。
\end{proof}

\endinput
